\section{Additional Recorded Experiments}
\label{app:f}\label{app:additional-experiments}
% Generated by figures/scripts/prepare_additional_experiments.py; do not edit.
\newcommand{\AdditionalHiddenObserved}{0.0611}
\newcommand{\AdditionalHiddenFull}{1.0000}
\newcommand{\AdditionalDiabetesGamma}{72.08}
\newcommand{\AdditionalDiabetesSelected}{0.6325}
\newcommand{\SupplementFiniteChoiceTable}{%
\begin{table}[!htbp]
\centering\small\setlength{\tabcolsep}{4pt}
\begin{tabular}{@{}lrcccc@{}}
\toprule
Rule & $f(B)$ & Observed CVaR & Full CVaR & Upper endpoint & Width\\
\midrule
Observed labels & .25 & 0.6447 & 0.7845 & 0.8877 & 0.2844\\
IPW & .25 & 0.6447 & 0.7845 & 0.8877 & 0.2844\\
Upper envelope & .65 & 0.6686 & 0.6582 & 0.8628 & 0.2286\\
Full-label oracle & .65 & 0.6686 & 0.6582 & 0.8628 & 0.2286\\
\bottomrule
\end{tabular}
\caption{Finite-candidate selection over 250 samples of size 1,400. Values average exact population evaluations of the chosen predictor; each rule selected the same candidate in every replication. The upper endpoint is a population sensitivity envelope, without statistical enlargement.}
\label{tab:additional-selection}
\end{table}
}
\newcommand{\SupplementFixedAuditTable}{%
\begin{table}[!htbp]
\centering\small\setlength{\tabcolsep}{4pt}
\begin{tabular}{@{}rccccc@{}}
\toprule
Labels & TIV & Random & Density & Uncertainty & Predicted loss\\
\midrule
\multicolumn{6}{l}{\emph{Fixed-threshold gap removed} (\%)}\\
600 & 42.4 & 30.6 & 40.7 & 15.8 & 15.8\\
1,200 & 76.2 & 61.6 & 67.8 & 47.1 & 47.1\\
1,800 & 100.0 & 91.4 & 95.5 & 88.9 & 88.9\\
\midrule
\multicolumn{6}{l}{\emph{Separately optimized empirical width removed} (\%)}\\
600 & 29.4 & 30.3 & 28.7 & 21.3 & 21.3\\
1,200 & 72.0 & 61.0 & 62.5 & 59.0 & 59.0\\
1,800 & 100.0 & 91.1 & 98.2 & 87.1 & 87.1\\
\bottomrule
\end{tabular}
\caption{One synthetic pool of 7,000 records, fixed predictor, and unit cost per hidden label. Reductions are normalized by that pool's initial gap or initial optimized width. These two targets need not rank policies identically; at 600 labels random revelation removes slightly more optimized width than TIV.}
\label{tab:additional-audit}
\end{table}
}
\newcommand{\SupplementReferenceForecastTable}{%
\begin{table}[!htbp]
\centering\small\setlength{\tabcolsep}{4pt}
\begin{tabular}{@{}lccc@{}}
\toprule
Specified forecast & Observed CVaR & Full CVaR & Upper endpoint\\
\midrule
\multicolumn{4}{l}{\emph{Synthetic populations: mean over 120 worlds}}\\
Selected-label probability $p_1$ & 0.8052 & 0.8703 & 0.8856\\
Sensitivity upper probability $q_+(p_1)$ & 0.8250 & 0.8102 & 0.8198\\
Full-label probability $p_{\rm full}$ & 0.7972 & 0.7464 & 0.8287\\
\midrule
\multicolumn{4}{l}{\emph{Diabetes-derived finite population}}\\
Selected-label probability $p_1$ & 0.2665 & 0.9252 & 0.9320\\
Sensitivity upper probability $q_+(p_1)$ & 0.8710 & 0.8693 & 0.8732\\
Full-label probability $p_{\rm full}$ & 0.5316 & 0.6760 & 0.7505\\
\bottomrule
\end{tabular}
\caption{Three specified reference forecasts, evaluated with absolute loss and $\alpha=.85$. The CSV role names ``Observed-selected,'' ``Certificate,'' and ``Oracle'' identify these formulas; they do not denote optimization over a candidate class. Synthetic entries use $\Gamma=12$; diabetes entries use its oracle-chosen sensitivity level.}
\label{tab:additional-forecasts}
\end{table}
}
\newcommand{\SupplementGridBridgeTable}{%
\begin{table}[!htbp]
\centering\small\setlength{\tabcolsep}{4pt}
\begin{tabular}{@{}rcccc@{}}
\toprule
 & \multicolumn{2}{c}{Synthetic median} & \multicolumn{2}{c}{Diabetes population}\\
Full-cost budget & Grid width & $W_{f,0}$ & Grid width & $W_{f,0}$\\
\midrule
0\% & 0.2562 & 1.3309 & 0.7295 & 1.1502\\
25\% & 0.0197 & 0.5393 & 0.2057 & 0.5857\\
50\% & 0.0080 & 0.1989 & 0.0084 & 0.1484\\
75\% & 0.0024 & 0.0599 & 0.0002 & 0.0025\\
100\% & 0.0000 & 0.0000 & 0.0000 & 0.0000\\
\bottomrule
\end{tabular}
\caption{Width connection along the fixed-$\eta=.5$ TIV-per-cost allocation. Each endpoint is minimized on the recorded 61-point cutoff grid. Synthetic columns are separate medians across 120 worlds; diabetes has one finite population. The sufficient bound is $W_{f,0}$, which is not itself a CVaR risk and can exceed one.}
\label{tab:additional-bridge}
\end{table}
}
\newcommand{\SupplementDiabetesCostTable}{%
\begin{table}[!htbp]
\centering\small\setlength{\tabcolsep}{4pt}
\begin{tabular}{@{}lr@{}}
\toprule
Allocation rule & First recorded cost to 25\% gap\\
\midrule
TIV/cost & 42\%\\
Fractional-knapsack frontier & 42\%\\
Random & 76\%\\
High density & 76\%\\
Lowest observation probability & 54\%\\
Predicted absolute loss & 44\%\\
Uncertainty & 44\%\\
Sensitivity width & 44\%\\
Excess-loss contrast & 86\%\\
\bottomrule
\end{tabular}
\caption{Diabetes-derived population: first hit on the 2\%-spaced full-cost grid. Random is the mean remaining gap over 100 random cell orderings, so its entry is the first crossing of that mean curve, not the mean of 100 crossing costs. Acquisition is divisible at the cell level.}
\label{tab:additional-diabetes-cost}
\end{table}
}


This appendix collects saved experiments that complement the studies in
Appendix~\ref{app:experiments}: simple observation mechanisms, finite-candidate
learning, additional diagnostics for the synthetic revelation frontier, and a
finite-population illustration constructed from diabetes data. They answer
questions about how the objectives behave under controlled conditions. Their
population quantities and empirical trajectories carry no additional
finite-sample confidence guarantee. No new simulations or model fits were run
when preparing these tables.

\subsection{Observation mechanisms, candidate learning, and fixed audits}
\label{app:additional-mechanisms}

\paragraph{Low observed risk can hide high population risk.}
The two-region example has masses $(.82,.18)$, deterministic outcomes $(0,1)$,
and observation probabilities $(.95,.04)$. Use absolute loss and $\alpha=.85$.
For the constant-zero predictor, the saved observed-label CVaR is
\AdditionalHiddenObserved{}, whereas population CVaR is
\AdditionalHiddenFull{}. The region-specific predictor $(0,1)$ has zero loss
under both populations. Thus this example demonstrates a large discrepancy
between observed and population risk; it does not exhibit a pairwise ranking
reversal. The saved envelope calculation uses $\Gamma=50$ and clips boundary
reference probabilities to $[10^{-8},1-10^{-8}]$, producing negligible nonzero
envelope values for the zero-loss reference.

\paragraph{Finite-candidate selection with known auxiliary functions.}
Use the five-region population specified under ``Population sweeps'' in
Appendix~\ref{app:finite-sample-experiments}. Outcomes are sampled from its
specified Bernoulli probabilities; observation is sampled conditional on the
outcome. The adverse-label observation probability in region $B$ is $.03$,
compared with $.60$ for the other label, giving an odds-ratio sensitivity shift
of $48.5$. Fix $\Gamma=50$, $\alpha=.85$, absolute loss, and
\[
 \mathcal F_0=\{(.05,b,.10,.40,.50):b\in\{.05,.25,.65,.85\}\}.
\]
Each of 250 independent training samples contains 1,400 records. Four rules
minimize, over $\mathcal F_0$, empirical full-label CVaR, empirical observed-label
CVaR, normalized inverse-propensity-weighted observed-label CVaR, or the empirical
upper sensitivity objective. The propensity rule uses the known marginal
$\pi(X)$; the upper objective uses known $p_1$ and combines observed excess loss
with the bounded hidden contribution. It does not estimate auxiliary models or
add the locally robust residual correction. All cutoff searches use exact loss
breakpoints. The full-label rule is an oracle benchmark unavailable to a learner
with selective labels.

Table~\ref{tab:additional-selection} evaluates each chosen predictor under the
exact population law. Each rule chooses the same candidate in every recorded
replication, so the table does not attach a sampling-error bar to its repeated
population evaluation. The difference between observed-label and
inverse-propensity learning, on one hand, and upper-envelope learning on the
other illustrates the consequence of outcome-dependent observation in this
particular candidate class. It is not a universal comparison between learning
algorithms.

\SupplementFiniteChoiceTable

\paragraph{A single fixed-predictor audit pool.}
A further sample of 7,000 records from the same population supplies one audit
pool. Fix the population observed-label-preferred predictor
$f=(.05,.25,.10,.40,.50)$ and the lower-envelope minimizing cutoff $\eta=.5$.
Only originally hidden outcomes can be acquired, at unit cost. The five static
rankings are per-hidden-label TIV $s_\Gamma|d_{f,\eta}|/\beta$, independent
uniform random scores, region mass, uncertainty $-|f-.5|$, and expected absolute
loss at the midpoint $(q_-+q_+)/2$ of the sensitivity interval. The implementation
sorts each score in decreasing order using NumPy's default \texttt{argsort}; it
does not impose a separate stable tie rule. The 16 recorded budgets range from
zero to 1,800 labels in increments of 120. The predictor and rankings remain
fixed throughout.

Revealed outcomes enter the known empirical contribution. The remaining hidden
outcomes retain the original sensitivity bounds. Table~\ref{tab:additional-audit}
separately reports the reduction of the fixed-cutoff gap and of the empirical
interval whose endpoints minimize their cutoffs independently. Both reductions
are normalized by their respective initial value in this one pool. The two
objectives need not favor the same allocation at every budget: random revelation
removes slightly more optimized width at 600 labels, although TIV removes more
fixed-cutoff ambiguity. Prediction uncertainty is also distinct from label-sensitive
tail ambiguity: the region with $f=.5$ has identical absolute loss under both
labels, hence zero TIV at every cutoff. These are mechanism observations from
one realized pool, without replication-based inference or a claim of adaptive
learning optimality.

\SupplementFixedAuditTable
\FloatBarrier

\subsection{Additional diagnostics for the synthetic frontier}
\label{app:additional-frontier}

Use the same 120 synthetic worlds, each with 120 cells, described in
Appendix~\ref{app:frontier-experiments}. The loss is absolute, $\alpha=.85$, and
$\Gamma=12$. The additional saved outputs compare three specified forecasts:
$p_1$, the upper sensitivity probability $q_+(p_1)$, and
$p_{\rm full}=\pi p_1+(1-\pi)p_0$. The CSV names for these rows are
``Observed-selected,'' ``Certificate,'' and ``Oracle.'' In this experiment they
are formulas fixed by the generator, not the results of minimizing empirical
risk or an upper endpoint over a candidate class. The full-label forecast uses
information unavailable from the observed law. Table~\ref{tab:additional-forecasts}
reports mean population evaluations over the 120 worlds, alongside the
single finite population introduced below. In particular, the $q_+(p_1)$ row
is not evidence that the risk-minimizing forecast has been found.

\paragraph{From a common cutoff to separately minimized endpoints.}
For each world, fix $f=p_1$ and acquire along the TIV-per-cost ordering at
$\eta=.5$. Cost is $\sum_j\mu_jc_j\lambda_j$, with a fractional terminal cell;
budget fractions divide by the full population cost $\sum_j\mu_jc_j$.
The known audited contribution is evaluated with the generative $p_0$, while the
remaining outcomes retain their sensitivity bounds. This isolates the retained
population objective from estimation of the newly audited contribution.

The saved diagnostic evaluates 21 cost fractions from zero to one. Its endpoint
search uses the grid
$\mathcal G=\{0,1/60,\ldots,1\}$, giving
\[
 I^{\mathcal G}_{\Gamma,\lambda}(f)
 =\min_{\eta\in\mathcal G}U_{f,\eta}(\lambda)
  -\min_{\eta\in\mathcal G}L_{f,\eta}(\lambda).
\]
This recorded \emph{grid-minimized width} is distinct from the exact
breakpoint-minimized width in Appendix~\ref{app:optimized-width}. The companion
bound maximizes the common-cutoff gap on the same grid. Because binary
absolute-loss contrasts satisfy $|d_{f,\eta}|\leq|d_{f,0}|$ and $0\in\mathcal G$,
that maximum equals $W_{f,0}(\lambda)$. All 2,520 saved world--budget rows satisfy
$0\leq I^{\mathcal G}_{\Gamma,\lambda}(f)\leq W_{f,0}(\lambda)$ to numerical
precision. Table~\ref{tab:additional-bridge} gives selected budgets. This check
illustrates the sufficient-width connection; it neither makes the $.5$-cutoff
allocation optimal for optimized width nor turns the finite grid into exact
threshold minimization. The gap bound can exceed one because it is not itself
a bounded-loss CVaR.

\subsection{A diabetes-derived finite-population illustration}
\label{app:additional-diabetes}

\paragraph{Measured outcomes and imposed observation.}
The construction uses all 442 records in the diabetes dataset distributed with
scikit-learn \citep{pedregosa2011scikit}. It binarizes the measured disease-progression
outcome at its empirical 75th percentile, assigning $Y=1$ at or above that
cutoff. A standardized logistic regression, fitted to all these outcomes with a
5,000-iteration limit and the recorded software defaults, supplies scores $r_i$.
Fourteen empirical-quantile score cells define the finite covariate population.
The score is used to construct strata, not evaluated as a held-out predictive
model. In particular, the construction has access to outcomes subsequently
called hidden.

For each record, impose the observation probability
\[
 s_i=\operatorname{clip}_{[.03,.95]}
 \{.85-.25r_i-.55Y_i-.20Y_i r_i+\epsilon_i\},
 \qquad \epsilon_i\sim N(0,.01^2).
\]
The saved construction integrates these probabilities; it does not draw a new
binary observation mask. For a cell $j$ with $n_j$ records, set
\[
 \mu_j=\frac{n_j}{442},\qquad
 \pi_j=\frac{1}{n_j}\sum_{i\in j}s_i,\qquad
 p_{1j}=\frac{\sum_{i\in j}Y_i s_i}{\sum_{i\in j}s_i},\qquad
 p_{0j}=\frac{\sum_{i\in j}Y_i(1-s_i)}{\sum_{i\in j}(1-s_i)}.
\]
Before evaluation, the implementation clips $p_{1j}$ and $p_{0j}$ to
$[10^{-4},1-10^{-4}]$ and uses the resulting Bernoulli mixture
$p_{{\rm full},j}=\pi_jp_{1j}+(1-\pi_j)p_{0j}$. Consequently, the reported full
risks refer to this smoothed finite-cell law. The measured progression outcomes
and covariates define the construction; the observation channel and acquisition
costs are imposed. The population observation share is
\AdditionalDiabetesSelected{}.

\paragraph{Sensitivity, loss, and costs.}
Use absolute loss and $\alpha=.85$. The sensitivity level is chosen with oracle
knowledge of both cell kernels:
\[
 \Gamma=\max\!\left\{2,\;1.10\exp\!\left(
 \max_j|\operatorname{logit}p_{0j}-\operatorname{logit}p_{1j}|\right)\right\}
 =\AdditionalDiabetesGamma{}\quad\text{(rounded)}.
\]
This construction ensures compatibility for the known cell law. It is not a
calibration procedure justified using observed labels alone, and it supplies no
simultaneous or data-dependent-$\Gamma$ confidence statement. With zero-based
cell indices, the imposed population cost density is
$c_j=.65+1.15\bar r_j+.20\mathbf1\{j\bmod3=0\}$, where $\bar r_j$ is the
cell's mean fitted score.

The three reference forecasts and their population risk evaluations follow
Appendix~\ref{app:additional-frontier}. For the acquisition comparison fix
$f=p_1$ and $\eta=.5$, use divisible cell allocations, and evaluate 51 fractions
of full cost in increments of $.02$. The acquisition rules are the seven
baselines and TIV-per-cost/frontier ordering described in
Appendix~\ref{app:frontier-experiments}. Random acquisition averages the remaining
gap over 100 independent uniform-random cell orderings. The 21-budget width
diagnostic uses the same 61-point cutoff grid as the synthetic study.

\SupplementReferenceForecastTable
\SupplementGridBridgeTable
\FloatBarrier

Table~\ref{tab:additional-diabetes-cost} reports the first recorded cost fraction
at which each curve reaches 25\% of its initial fixed-cutoff gap. For random
acquisition, this is the crossing of the mean gap curve, rather than the mean of
100 individual crossing costs. TIV and the exact fractional-knapsack frontier
coincide by their common ordering. The data illustrate how observation and costs
shape a retained population certificate when all relevant quantities are known;
they do not establish clinical effectiveness, generalization to a new patient
population, or finite-sample coverage. They also do not convert passive revelation
into identification of outcomes under untaken actions.

\SupplementDiabetesCostTable
\FloatBarrier

\subsection{Saved-output provenance and reproduction order}
\label{app:additional-provenance}

The numerical sources are the \texttt{experiment1}, \texttt{experiment2}, and
\texttt{experiment3} CSVs, and the \texttt{experiment6} and
\texttt{experiment7\_semisynthetic} CSVs under \path{experiments/results/}.
The mechanism examples use the full-sequence random stream with master seed
20260503. The current frontier and diabetes CSVs instead match
\path{run_summary_main_figure.json}, generated by running
\path{experiments/run_tail_observability_experiments.py} with
\texttt{--main-figure}. The full-run summary's older frontier output uses a different random-stream
position. The archived \path{run_summary_semisynthetic.json} instead records an
earlier 569-record malignancy example, with malignancy-rate fields; it is not
the 442-record diabetes construction in this appendix. Neither archived output
is substituted for the current CSVs.
Reproductions should use separate output copies or follow the documented run
order so that one entry point does not silently overwrite another run's saved
results.

The standard-library extraction script
\path{figures/scripts/prepare_additional_experiments.py} regenerates these tables
from the saved CSVs. Its \texttt{--check} mode verifies the generated tables,
source checksums, agreement with the matching run summary, repeated-sample
aggregation, TIV/frontier equality, and the recorded width inequalities.
\path{figures/supplement_additional_provenance.json} records those checks and the
source-to-entry-point mapping. The corrected lending sequential experiment,
original-cohort verification, and finite-sample endpoint study remain documented
in Appendix~\ref{app:experiments} and the accompanying reproduction material.
